% Appendix E, from exam/paper_b.md: Paper B, questions only. The model answers are a separate
% appendix, as they are a separate file, so the paper can be sat without them in view.

\chapter{Self-Assessment Paper B}
\label{exb:ch:paper}

\noindent\textbf{Time:} 4 hours. \textbf{Closed book.} A basic calculator is permitted.
\textbf{Total: 100 marks.}

\medskip
This paper exists to test your own skills and knowledge. It is not a qualification and it gates
nothing. It draws on all ten chapters and on both halves of the system, so \textbf{it is meant to be
taken once you have finished the book}, after \chapref{10} and the bring-up ladder. The model
answers are in Appendix~\ref{ansb:ch:answers}.

\section*{Rubric}
\textbf{The protocol is the single source of truth.} Where an implementation and the protocol in
Appendix~\ref{proto:ch} disagree, the implementation is wrong. An answer that argues from the
contract beats one that argues from code.

\textbf{Both languages are the subject, so you write both here.} VHDL is VHDL-93 using
\code{ieee.std_logic_1164}, and \code{ieee.numeric_std} where a conversion needs it; write the
\code{library} and \code{use} clauses out at least once, after which they are assumed. C++ that would
run on the ATmega328P is written in the book's AVR-portable subset: \code{<stdint.h>} and the bare
types, no \code{std::}, nested \code{namespace} blocks rather than \code{namespace driver::uart}, no
\code{[[nodiscard]]}, no \code{inline} variables. Host-only test code may use full modern C++, and a
question says so when it means it.

\textbf{Syntax is not what is being marked. The design is.} A missing semicolon, a forgotten
\code{is} or an absent \code{#include} costs nothing. A process that describes a latch where you
meant a wire, or a \code{readReg} that assembles its four bytes backwards, costs that part in full,
because that is a different system.

\textbf{The book's conventions are the ones expected}, and a design that breaks one should say why:
\begin{itemize}
  \item VHDL ports are declared all inputs first, then all outputs, and every \code{port map} and
    \code{generic map} is \textbf{positional}.
  \item An active-low signal is named with an \code{_n} suffix; a signal that has been through a
    two-flop synchronizer is named with an \code{_s2} suffix. Every submodule takes
    \code{reset_s2_n}; only \code{uart_top} takes the raw \code{reset_n}. Signals carrying SPI
    traffic between the two transport blocks take a \code{spi_} prefix.
  \item Resets are asserted \textbf{asynchronously}, so \code{reset_s2_n} belongs in the
    sensitivity list.
  \item C++ uses \code{camelCase}, a leading capital on type names, \code{my} on private members
    and \code{our} on private statics, brace initialization, \code{noexcept} on driver methods, and
    copy and move deleted on any class holding a reference member or a unique hardware resource.
\end{itemize}

\textbf{Where a question asks you to review a design}, naming the defect earns half the marks and
saying what it does to the running system earns the other half. ``This is wrong'' scores nothing.

\subsection*{Supplied constants and formulas}

\begin{booktable}{@{}p{12em}L@{}}
  \thead{Quantity} & \thead{Value} \\
  \midrule
  DE0-CV system clock & 50\,MHz, a 20\,ns period \\
  Arduino Nano clock (f\textsubscript{osc}) & 16\,MHz \\
  SPI \code{SCK} & f\textsubscript{osc}/16, which is 1\,MHz; mode 0, MSB first \\
  Receiver oversampling & 16 times \\
  Default frame & 8N1: 1 start bit, 8 data bits LSB first, 1 stop bit \\
  Baud divider & \code{BAUD_DIV = round( 50_000_000 / (16 * baud) )} \\
  SPI transaction & 5 bytes: 1 command byte, then 4 data bytes MSB first \\
  SPI command byte & bit 7 = \code{W} (1 write, 0 read), bits 6--4 = 0, bits 3--0 = register
    index \\
\end{booktable}

The register map itself is \textbf{not} supplied: Question~\ref{exb:q:5} asks you to write it out,
and Questions~\ref{exb:q:4} and~\ref{exb:q:6} consume it, which is what the follow-through rule is
for.

\subsection*{Marks}

\begin{narrowtable}{@{}lcccccccc@{}}
  \thead{Question} & 1 & 2 & 3 & 4 & 5 & 6 & 7 & 8 \\
  \midrule
  \thead{Marks} & 12 & 13 & 14 & 13 & 12 & 13 & 11 & 12 \\
\end{narrowtable}

% ------------------------------------------------------------------------------------------
\question{Building the top before anything it holds}{12}
\label{exb:q:1}

\qpart{a} Write the complete \code{uart_top} \textbf{entity}, with its \code{library} and \code{use}
clauses and the book's banner comment naming the inputs and outputs. Eight ports, all inputs first,
in the order the system testbench binds them.

Then name the \code{use} clause the file does \textbf{not} need in \chapref{1} and say why, name the
chapter that adds it and the single expression that needs it, and state why \code{uart_top} never
writes \code{use work.uart_def.all} at all. \worth{3}

\qpart{b} Here are the two provided transport blocks' port lists, in declaration order.

\begin{console}
spi_slave       1 clock     2 reset_s2_n  3 sclk       4 mosi      5 ss
                6 tx_data   7 miso        8 rx_data    9 rx_valid  10 ss_active

spi_reg_bridge  1 clock     2 reset_s2_n  3 ss_active  4 rx_data   5 rx_valid
                6 reg_rdata 7 tx_data     8 reg_addr   9 reg_wdata 10 reg_write
\end{console}

\noindent\code{spi_slave}'s ports 1--5 are \code{std_logic} inputs, 6 is an 8-bit input vector, 7 is
a \code{std_logic} output, 8 is an 8-bit output vector, and 9--10 are \code{std_logic} outputs.
\code{spi_reg_bridge}'s ports 1--3 are \code{std_logic} inputs, 4 is an 8-bit input vector, 5 is a
\code{std_logic} input, 6 is a 32-bit input vector, 7 is an 8-bit output vector, 8 is a 4-bit
output vector, 9 is a 32-bit output vector, and 10 is a \code{std_logic} output.

Declare \textbf{exactly} the internal signals these two instantiations and the provided
\code{reset_sync} require, with their types, following the book's naming conventions, and write the
three positional instantiations. Do not declare anything a later chapter will bring.

Then answer two things. \code{spi_slave} and \code{spi_reg_bridge} both have a port called
\code{rx_data}: state why those are not one bus and which way each of your signals runs. And state
why \code{reg_addr} is four bits wide when there are only seven registers. \worth{5}

\qpart{c} \chapref{1} ends with one placeholder assignment, because the register bank is four
chapters away.

Write it. Then state what an SPI \textbf{read} transaction returns while it is in place and why that
is harmless until \chapref{5}, and state what a read would return instead if the placeholder were
left out altogether --- being precise about what ``returns'' means for a signal nobody drives.
\worth{2}

\qpart{d} State what building \code{uart_top} before any block it instantiates fixes about every
block you write later.

Then: \code{baud_gen} arrives in \chapref{2}. State exactly what has to be true of its entity for
its instantiation in \code{uart_top} to bind correctly without you touching the top again, and name
the one thing about it that the book explicitly says is \emph{not} part of the contract. \worth{2}

% ------------------------------------------------------------------------------------------
\question{The generator, the divider, and the frame}{13}
\label{exb:q:2}

\qpart{a} Write \file{baud_gen.vhd} complete: the \code{library} and \code{use} clauses, the entity
with its four ports in order and their exact types, and the architecture. It produces a registered,
one-clock-wide \code{tick} every \code{div} clocks, and its reset is the book's.

Then justify your comparison. State what \code{counter >= div - 1} and \code{counter = div - 1}
each do when \code{div} is \code{1}, and what each does if \code{div} could ever be \code{0}, and
say which keeps the counter inside its declared range. \worth{4}

\qpart{b} The peripheral is to run at 9600 baud.

Give the exact (unrounded) divider, the value written to \code{BAUD_DIV}, the baud rate achieved,
and the error as a percentage with its sign.

Then state why the division is by \code{16 * baud} rather than by \code{baud}, name the module that
spends the sixteen ticks and say what it spends them on, and state what the transmitter does with
the same tick stream. \worth{3}

\qpart{c} \code{uart_tx} produces 8N1 only. Extend it with \code{parity_en} and \code{parity_odd}
inputs.

Give the new frame vector expression and its index range, say which index the parity bit occupies
and which the stop bit now does, and give the expression that computes the parity bit from
\code{data} for even parity and for odd.

Then a \code{two_stop} input is added as well. Name the only field of the frame whose length now
varies, give the three possible frame lengths in bits, and state what the longest frame costs in
transmitted bytes per second at 115200 baud compared with 8N1. \worth{3}

\qpart{d} A candidate's transmitter passes \code{uart_tx_tb}. Their \code{busy} is not the book's
concurrent line but a registered flag, set inside the clocked process on the edge \emph{after} the
state becomes \code{STATE_SEND}, so it rises one clock late.

State what \code{uart_tx_tb} checks about \code{busy} and what it does not. Then state what
\code{uart_top}'s TX feeder does during that one clock, what \code{uart_tx} does with what the
feeder sends, and what happens to the byte involved. Give the one-line fix and say which chapter's
testbench is the first that could have caught it. \worth{3}

% ------------------------------------------------------------------------------------------
\question{The asynchronous line, written out}{14}
\label{exb:q:3}

\qpart{a} Write \file{sync.vhd} complete: the \code{library} and \code{use} clauses, the entity with
its generic and its four ports in order, and the architecture. Two flip-flops per bit, the book's
reset, and the output taken from the right one.

Then answer two things. State how many rising edges \code{sync_out} takes to follow
\code{async_in}, and why neither one fewer nor one more is correct. And state what a
\code{COUNT = 8} instance does \textbf{not} guarantee about eight bits that change on the same edge,
why that is not a defect in \code{sync}, and one kind of signal you must therefore never put
through a wide instance. \worth{4}

\qpart{b} Write the \code{STATE_IDLE} and \code{STATE_START} branches of \code{uart_rx}'s state
machine as the book specifies them, including the one flip-flop of history, where in the process it
is updated, and why that position matters. Assume the surrounding \code{if baud_tick = '1'} and the
tick counter are already written.

Then a candidate replaces the falling-edge test with a level test, \code{if rx_s2 = '0' then}. State
what this does during a \textbf{break} --- the line held low for longer than one frame --- and then
state precisely what happens to the frame that is still in flight at the moment the break
\emph{ends}. Name the testbench case that catches it and say what its message tells you. \worth{4}

\qpart{c} Write the \code{STATE_STOP} branch, and the output defaults at the top of the clocked
process that go with it.

State what makes \code{valid} and \code{frame_err} exactly one clock wide, why the byte is dropped
on a framing error rather than delivered alongside a flag, and what the register bank does with each
of the two pulses. \worth{3}

\qpart{d} An earlier revision of \code{uart_rx_tb} sent \code{0xA5} in its clean-frame case and
\code{0x3C} elsewhere, while \code{uart_top_tb} sent \code{0x5A} through a \code{tx}-to-\code{rx}
loopback.

Show that a receiver which stores its data bits in the wrong order passed all three. Name the
property \code{0xA5}, \code{0x3C} and \code{0x5A} share, state how many of the 256 bytes have it,
and give the \textbf{second}, independent reason \code{uart_top_tb} could not have caught a
bit-order convention shared by \code{uart_tx} and \code{uart_rx}, however good its constant was.

\code{uart_rx_tb}'s clean-frame case sends \code{0x53} today. State the value \code{data_out}
carries under the faulty receiver, and hence what the testbench now reports. Then state which single
kind of bit-order mistake remains beyond the reach of every testbench in \file{hw/}, and why.
\worth{3}

% ------------------------------------------------------------------------------------------
\question{The bank, written out}{13}
\label{exb:q:4}

\qpart{a} Write the \code{ERROR_FLAGS} latch for the framing bit: its reset value, the condition
that sets it, the condition that clears it, and what happens when both occur on the same clock. Use
\code{uart_def}'s constants rather than numbers.

Then state why a single-cycle error pulse cannot be polled by software at all, and say what would
actually be observed if \code{STATUS} bit 2 were wired straight to the receiver's \code{frame_err}
output and a driver polled \code{STATUS} once per SPI transaction. \worth{4}

\qpart{b} Write the bank's read path: the combinational decode that drives \code{reg_rdata} from
\code{reg_addr}, covering all seven registers and the reserved indices, using \code{uart_def}'s
constants. Then write the two write \textbf{actions}, for \code{TX_DATA} and for \code{RX_POP}, as
the book's bank implements them.

State why the read is combinational rather than registered, what the \code{when others} branch must
produce and why getting it wrong would be a different kind of mistake here than in a purely
combinational decoder, and what the bank relies on the provided bridge to guarantee about
\code{reg_write}. \worth{4}

\qpart{c} Overrun is reserved in this course but not implemented. Design it.

Give the condition that detects an overrun, in terms of signals \code{uart_regs} already has. Give
the \code{ERROR_FLAGS} bit it sets and its position, and the clear path it shares. State the one
piece of information \code{uart_rx} would have to be given for it to detect the overrun itself
instead, and why handing it that information is a worse design.

Then give a physical cause for an overrun and a different physical cause for a framing error.
\worth{3}

\qpart{d} A candidate's bank stores \code{STATUS} in a 32-bit register, updated in the clocked
process whenever \code{rx_push}, \code{tx_pop} or \code{frame_err} asserts. It passes a test they
wrote themselves.

Give two distinct ways it disagrees with \code{uart_regs_tb}, and state the general rule the book
draws from it in one sentence. \worth{2}

% ------------------------------------------------------------------------------------------
\question{Writing the contracts}{12}
\label{exb:q:5}

\qpart{a} Write \file{include/driver/transport/interface.hpp} and
\file{include/driver/uart/interface.hpp} in full, in the book's AVR-portable style: the include
guard or \code{#pragma}, the includes, the namespaces, the destructors, and every pure virtual method
with its exact signature, parameters, return type and qualifiers.

\code{driver::transport::Interface} has three operations. \code{driver::uart::Interface} has six.
\worth{5}

\qpart{b} Write the \code{reg} and \code{status} nested namespaces of
\file{include/driver/uart/register_map.hpp}, with the correct names and values, in the book's style.

Then state why the constants are declared plain \code{constexpr} rather than
\code{inline constexpr}, and what ``internal linkage at namespace scope'' costs here. And state why
the map stores bit \textbf{positions} rather than masks, naming the file on the other side of the
wire that this decision keeps identical and the one-line idiom that forms a mask at the use site.
\worth{3}

\qpart{c} \code{driver::uart::Stub} holds \code{volatile bool& myStop} as a member.

State what that reference is for and who owns the flag it refers to. State what \code{read()} does
with it when the scripted RX buffer runs out, and why that behaviour exists at all given that a stub
has no reason to care how a caller's loop terminates.

Then state what holding a reference member forces about the class's copy constructor, move
constructor and default constructor, and why. \worth{2}

\qpart{d} \code{driver::uart::Stub}'s \code{status()} returns \code{0}, its \code{errorFlags()}
returns \code{0} and its \code{clearErrors()} does nothing, even though a real UART's \code{STATUS}
is the register the whole driver is built around.

State the general rule about what a test double must be faithful about and what it may ignore, and
apply it here: say what would have to change about the stub before it could be used to test
something new, and give an example of such a thing. \worth{2}

% ------------------------------------------------------------------------------------------
\question{Writing the protocol}{13}
\label{exb:q:6}

\qpart{a} Write \code{Uart::readReg(uint8_t addr) const} and
\code{Uart::writeReg(uint8_t addr, uint32_t value)} in full, in the book's style, using only
\code{myTransport}.

Both perform one five-byte transaction. Get the command byte, the framing and the byte order right,
and say in one line each what \code{readReg} does with the reply to the command byte and what
\code{writeReg} does with all four of its replies.

Then state what marking \code{readReg} \code{const} buys the class, and name the two public methods
that depend on it. \worth{5}

\qpart{b} Write \code{Uart::read(uint8_t& byte)}.

Then, for one successful call on a peripheral whose RX FIFO front byte is \code{0x41}, write out the
exact bytes a \code{driver::transport::Stub} would record, transaction by transaction, and give
\code{beginCalls()} and \code{endCalls()}.

Finally, write the scripting the test must perform before the call. The stub's helpers are
\code{injectRxByte(uint8_t)} and \code{injectRxWord(uint32_t)}, the latter queuing four bytes most
significant first. State how many bytes in total have to be queued for this one call and explain the
discrepancy between that number and the number of register values involved. \worth{4}

\qpart{c} Write \code{writeBlocking()} and \code{readBlocking()} from
\file{driver/uart/blocking.hpp}.

State why each is \code{inline}, why each takes \code{Interface&} rather than \code{Uart&}, and why
they are free functions in a separate header rather than methods on the interface. \worth{2}

\qpart{d} A candidate implements \code{readReg} by calling the driver's own \code{write()} to send
each byte, reasoning that \code{write()} is ``the method that sends a byte''.

State what happens when \code{status()} is called on this driver, in terms of the call sequence, and
name the failure the ATmega328P actually exhibits. Then state the one-sentence rule that keeps the
two apart. \worth{2}

% ------------------------------------------------------------------------------------------
\question{Writing the transport}{11}
\label{exb:q:7}

\qpart{a} Write \code{driver::transport::AvrSpi} in full: the constructor, the destructor and the
three seam methods, in the book's style, reaching the registers through the platform header's
\code{SCK}, \code{MOSI}, \code{MISO} and \code{SS} bit positions and the \code{DDRB}, \code{PORTB},
\code{SPCR}, \code{SPSR} and \code{SPDR} registers.

The class configures the ATmega328P as an SPI master at f\textsubscript{osc}/16, mode 0, MSB first.

Then state exactly what the destructor must clear, what it must leave alone and why, and state why
the constructor writes \code{SPCR = ...} rather than \code{SPCR |= ...}. \worth{5}

\qpart{b} A candidate ships \code{~AvrSpi() noexcept override = default;}, on the grounds that the
class has no members to release and the program never destroys it anyway.

Name the principle broken and list what the constructor took that is now never given back. Then
give one concrete way this harms code that runs afterwards, and one reason ``the program never
destroys it'' is not a safe assumption in this codebase in particular. \worth{3}

\qpart{c} The SPI transport needs no \code{F_CPU}; the debug logging over the ATmega's own USART
does.

State what sets the SPI \code{SCK} rate and what sets the USART's baud rate, and use that to explain
the difference.

Then: the Nano's fuses actually run it at 8\,MHz, but the build still defines
\code{F_CPU=16000000UL}. State what happens to \code{SCK}, in figures, and whether the SPI link still
works. State what happens to the debug USART, in figures, and what appears in the terminal. Say
which of the two failures you would notice first and why that is unfortunate. \worth{3}

% ------------------------------------------------------------------------------------------
\question{The application, and the plan}{12}
\label{exb:q:8}

\qpart{a} Write \code{app::EchoNode} in full --- \file{include/app/echo_node.hpp} and
\file{source/app/echo_node.cpp} --- in the book's AVR-portable style. It implements
\code{app::Interface}, whose single operation is \code{void run(const volatile bool& stop) noexcept}.

Include the member, the constructor, the deleted operations and the loop body. State why the
receive is the non-blocking \code{read()} while the echo is \code{writeBlocking()}, and why
\code{stop} is a \code{volatile bool} passed by \code{const} reference rather than an atomic or a
return value. \worth{4}

\qpart{b} Describe the host test for it precisely enough to write: which stub it runs over, what it
queues, how the loop is made to terminate, and what it asserts and in what order. Host test code may
use full modern C++.

Then state what an additional case that queues \textbf{nothing} proves, name the implementation
mistake it isolates, and say whether the three-byte case can catch that mistake too. \worth{3}

\qpart{c} Give the five rungs of the bring-up ladder, in order, and for each name the single layer
it adds that the previous rung did not exercise.

State why the order is forced rather than a matter of taste, naming the register that settles it
and the reason it can only be reached one way.

Then a colleague proposes flashing \code{app::EchoNode} first: ``if it echoes, everything works, and
if it does not we start debugging.'' Give two reasons this is a worse plan \textbf{even when it
succeeds}. \worth{3}

\qpart{d} A candidate writes \code{EchoNode} to hold \code{driver::uart::Uart& myUart} rather than
\code{driver::uart::Interface& myUart}. Everything compiles, the bench works, and the class is
otherwise identical.

State what is now impossible, what has to be built before the class can be tested at all, and give
the one-line fix. Then name the design decision of \chapref{6} that this makes concrete. \worth{2}
